% Introduction to the EU: Exam Strategy & Timeline (Day Before)
% Compile with: pdflatex exam-strategy.tex

\documentclass[8pt,landscape,a4paper]{article}
\usepackage[utf8]{inputenc}
\usepackage[T1]{fontenc}
\usepackage[left=0.8cm,right=2.6cm,top=0.8cm,bottom=0.8cm,marginparwidth=2cm,marginparsep=0.3em]{geometry}
\usepackage{multicol}
\usepackage{marginnote}
\usepackage{enumitem}
\setlist{nosep,leftmargin=*}
\usepackage{booktabs}
\usepackage{parskip}
\usepackage{microtype}
\usepackage{xcolor}
\usepackage{fancyhdr}
\definecolor{eublue}{RGB}{0,51,153}
\definecolor{charcoal}{RGB}{51,51,51}
\definecolor{softgrey}{RGB}{240,240,240}
\definecolor{notegrey}{RGB}{90,90,90}
\newcommand{\handnote}[1]{\marginnote{\raggedright\footnotesize\color{notegrey}\textit{#1}\vspace{0.4em}}}
\setlength{\parskip}{0.2em}
\setlength{\columnsep}{0.5em}
\pagestyle{fancy}
\fancyhf{}
\fancyfoot[L]{\raisebox{0.8em}{\small\color{notegrey} Intro EU}}
\fancyfoot[R]{\raisebox{0.8em}{\small\color{notegrey}\thepage}}
\fancypagestyle{plain}{\fancyhf{}\fancyfoot[L]{\raisebox{0.8em}{\small\color{notegrey} Intro EU}}\fancyfoot[R]{\raisebox{0.8em}{\small\color{notegrey}\thepage}}}
\raggedright

\begin{document}
\begin{center}
\vspace{0.2em}
{\normalsize\bfseries\color{eublue} Introduction to the EU Exam Strategy (Day Before)}\\[0.2em]
\footnotesize\textcolor{notegrey}{Use this with your main midterm guide. Goal: overview of the process + how to answer like your professor.}\\[0.35em]
\noindent\colorbox{softgrey}{\parbox{0.92\textwidth}{\centering\footnotesize\color{charcoal}
\textbf{Key:} \textbf{Strategy} \quad \textcolor{eublue}{\textbf{Timeline}} \quad \textit{Professor tips}
\vspace{0.25em}}}
\end{center}
\vspace{-0.4em}

\begin{multicols}{2}

\section*{\color{eublue}1. What the exam is (and isn’t)}\handnote{Not tricky, but you must show overview.}
\begin{itemize}[leftmargin=*]
\item \textbf{Not a puzzle:} Professor repeats that the exam is \emph{not} tricky or “genius level” but you must have studied.
\item \textbf{Task:} Read short texts, identify main idea, argument, key expression; match to course themes and periods.
\item \textbf{You need:} One clear mental map of European integration and institutions, not every detail.
\item \textbf{Time:} Around 90 minutes “a lot of time”. Use it to read, plan, and then write structured answers.
\end{itemize}

\section*{\color{eublue}2. Timeline of integration (5 stages)}\handnote{Turn this into a picture before the exam.}
\begin{itemize}[leftmargin=*]
\item \textbf{1945 to 1958} Foundations after WWII. Marshall Plan, Council of Europe, NATO, ECSC, early failed defence plans.
\item \textbf{1958 to 1973} \textbf{Treaty of Rome to first oil crisis.} EEC + Euratom, customs union, common market, CAP, Luxembourg Compromise, first enlargement.
\item \textbf{1973 to 1992} Crises and adaptation. Oil shock, stagflation, Southern enlargement, Single European Act, path to Maastricht.
\item \textbf{1992 to present} Maastricht to now. Euro, pillars, Schengen, Lisbon, enlargements, new policies (Green Deal, etc.).
\item \textbf{2008 crisis} Financial and euro crisis as a new phase inside this long period.
\end{itemize}

\subsection*{\color{eublue}How to use the timeline}
\begin{itemize}[leftmargin=*]
\item Place every agreement, institution change, and reading into one of these stages.
\item When you see a text in the exam, ask: \emph{Which stage? Which challenge? Which response?}
\item Use one “middle” event per stage as a memory hook (e.g.\ Rome, oil crisis, SEA, Maastricht, 2008).
\end{itemize}

\section*{\color{eublue}3. Key agreements to recognise}\handnote{Enough to identify them and say why they matter.}
\begin{itemize}[leftmargin=*]
\item \textbf{Treaty of Rome} (1957) Creates EEC + Euratom. Customs union and common market project.
\item \textbf{Common Agricultural Policy (CAP)} Early common policy; French priority; shows depth of economic integration.
\item \textbf{Luxembourg Compromise} (1966) ``Vital national interest'' veto; freezes supranational push for years.
\item \textbf{Single European Act} (1986) Single market; more qualified majority voting.
\item \textbf{Maastricht} (1992) EU, euro path, citizenship, pillars; big step beyond purely economic logic.
\end{itemize}

\section*{\color{eublue}4. Dimensions of integration}\handnote{Economic most complete; others have limits.}
\begin{itemize}[leftmargin=*]
\item \textbf{Economic:} Most complete customs union, single market, euro, common policies (CAP, competition).
\item \textbf{Political:} Increasing but incomplete institutions with real powers, but strong role for states and vetoes.
\item \textbf{Military / foreign policy:} Clear limits NATO, sovereignty, different national traditions.
\item In answers, show you know which dimension the text speaks to and how far integration goes there.
\end{itemize}

\section*{\color{eublue}5. Institutions overview (video recap)}\handnote{Use this with the institutions section of midterm.tex.}
\begin{itemize}[leftmargin=*]
\item \textbf{Commission} Engine room; proposes laws; guardian of treaties; one commissioner per portfolio.
\item \textbf{Parliament} Only directly elected body; represents citizens; co-legislator in most areas.
\item \textbf{Council of the EU} Ministers by policy area; co-legislator with Parliament.
\item \textbf{European Council} Heads of state and government; sets big political direction (not same as Council of the EU).
\item \textbf{Other bodies} Court of Justice, ECB, Court of Auditors, agencies.
\end{itemize}

\section*{\color{eublue}6. How to answer like the professor}\handnote{Copy his style.}
\begin{itemize}[leftmargin=*]
\item \textbf{Name the concept} first: e.g.\ “This is a \emph{low cost strategy based on task specialisation and standardisation of components}.”
\item \textbf{Define in one line}: short, textbook-style phrase.
\item \textbf{Two-three pieces of evidence}: from the text (dates, quotes, examples).
\item \textbf{Conclude}: link back to the question (e.g.\ “Therefore, Dacia follows a cost leadership strategy.”).
\item Use the same technical words he uses in model answers (heterogeneity, information asymmetry, etc.).
\end{itemize}

\section*{\color{eublue}7. Last-day checklist}\handnote{Run through this tonight.}
\begin{itemize}[leftmargin=*]
\item Can you sketch the 5-stage timeline (1945 to 1958, 1958 to 1973, 1973 to 1992, 1992 to present, 2008 crisis)?
\item Can you explain Rome, CAP, Luxembourg Compromise, SEA, Maastricht in 2 to 3 lines each?
\item Can you describe Commission, Parliament, Council of EU, European Council without looking?
\item Can you summarise one or two readings for each big period (origins, Southern enlargement, institutions, policies)?
\item Do you know how you will structure each answer (name-define-evidence-conclude)?
\end{itemize}

\end{multicols}

\end{document}

